\begin{proof}[Proof of \cref{obj:prop:oracle-regime-reduction}]
\begin{enumerate}
\item The well-formed-constants and baseline-slack hypotheses assumed here — \(\alpha>0\),
\(\beta>0\), \(s>0\), \(M>0\), \(c_0>0\), \(t_0\in(0,1)\), \(\varepsilon_0\in(0,1/2)\),
\([t_0-\varepsilon_0,t_0+\varepsilon_0]\subset(0,1)\), and \cref{obj:ass:baseline-submodel-slack}
for \((d,\beta,s,M,c_0,\varepsilon_0,t_0)\) — are exactly the hypotheses of
\cref{obj:thm:sharp-pointwise-lower-bound}. (The smooth-covariate condition \(d\le4s\) is not used
here; it enters only in Step 3.) That theorem therefore supplies a constant \(c>0\) and a threshold
\(n_1\) such that
\[
c\, n^{-2\alpha/(2\alpha+1)}
\le
R_n\!\left(\mathcal P_{\alpha,\beta,s}(M,c_0,\varepsilon_0,t_0),t_0\right)
\qquad\text{for every }n\ge n_1 .
\]
Fix this \(c\); it is the constant asserted in the statement. % lean: sharp_pointwise_lower_bound

\item Set \(n_0=\max\{n_1,1\}\). For every \(n\ge n_0\) both the display of Step 1 and the
elementary sample-size condition \(n\ge1\) are available; this is just the intersection of two
eventual conditions, and the conclusion of the proposition is asserted only for such \(n\).
% lean: Filter.eventually_ge_atTop

\item Fix \(n\ge n_0\). The hypotheses \(\alpha>0\), \(s>0\), the smooth-covariate condition
\(d\le4s\), and \(n\ge1\) are exactly those of \cref{obj:lem:rho-oracle-regime-algebra}, which
therefore gives the first asserted conclusion,
\[
\rho_n
=
n^{-2\alpha/(2\alpha+1)} .
\]
% lean: rho_oracle_regime_algebra

\item Substituting this identity into the display of Step 1 — that is, rewriting
\(n^{-2\alpha/(2\alpha+1)}\) as \(\rho_n\) on the left-hand side — yields the second asserted
conclusion,
\[
c\,\rho_n
\le
R_n\!\left(\mathcal P_{\alpha,\beta,s}(M,c_0,\varepsilon_0,t_0),t_0\right),
\]
valid for every \(n\ge n_0\). Together with \(c>0\) and the identity of Step 3, this is the claim.
% lean: oracle_regime_reduction
\end{enumerate}
\end{proof}
